\chapter{Bibliography Style}\label{app:bibliography-style}

The rules below explain how to properly format the bibliography of your thesis.
This template uses the standard Bib\TeX{} style \lstinline|plainurl|, a variant
of \lstinline|plain| that prints DOIs as links
(see \lstinline|\bibliographystyle| in \lstinline|main.tex|).
All entries of this chapter are contained in \lstinline|bibliography.bib|
and are cited here, so you can compare the source of each entry with the
resulting output in the bibliography at the end of this thesis.

%-----------------------------------------------------------------------
\section{Authors}\label{sec:bib-authors}
%-----------------------------------------------------------------------

\begin{itemize}
    \item Full first names, all authors, never ``et al.'' inside the entry:\\
          \lstinline|author = {Lommen, Nils and Meyer, {\'E}l{\'e}anore and Giesl, J{\"u}rgen}|
    \item Accents in \LaTeX{} form:\\
          \lstinline|Giesl, J{\"u}rgen|,
          \lstinline|Meyer, {\'E}l{\'e}anore|, \lstinline|Thiemann, Ren{\'e}|.
    \item Name particles lowercase and braced, so they stay with the surname:\\
          \lstinline|{von Essen}, Christian|.
    \item Hyphenated surnames braced: \lstinline|{Schneider-Kamp}, Peter|.
\end{itemize}

%-----------------------------------------------------------------------
\section{Titles}\label{sec:bib-titles}
%-----------------------------------------------------------------------

\begin{itemize}
    \item Title Case for every main word, protected with braces:
\begin{lstlisting}
{{Proving}} {{Almost-Sure}} {{Innermost}} {{Termination}} of
{{Probabilistic}} {{Term}} {{Rewriting}} {{Using}} {{Dependency}} {{Pairs}}
\end{lstlisting}
    \item Tool names in sans-serif via the macro \lstinline|\tool|, which is
          defined in \lstinline|verifyThesisStyle.sty|, so that every tool name
          looks the same in the text and in the bibliography:
          \lstinline|{{\tool{AProVE}}}| yields \aprove{}.
          The double braces protect the macro from Bib\TeX's case changes:
\begin{lstlisting}
{{\tool{KoAT}}}: {{Automatic Complexity}} and {{Termination Analysis}}
of {{Integer Programs}}
\end{lstlisting}
    \item Mathematical symbols in titles as math: \lstinline|Non-$E$-Overlapping|.
\end{itemize}

%-----------------------------------------------------------------------
\section{Venues}\label{sec:bib-venues}
%-----------------------------------------------------------------------

\paragraph{Conferences.}
Use one fixed pattern, always the same for the same series:
\lstinline|Proc.\ ACRONYM~'YY|, then series with volume, pages, and year.
No ``Proceedings of the 30th International Conference on \ldots'', no ordinal,
no city, no publisher.
The following two entries are rendered as \cite{kassinggiesl2023iast} and
\cite{ahrens2025WeightedRewritingSemiring}.

\begin{lstlisting}
@inproceedings{kassinggiesl2023iast,
  author    = {Kassing, Jan-Christoph and Giesl, J{\"u}rgen},
  title     = {{{Proving}} {{Almost-Sure}} {{Innermost}} {{Termination}} of
               {{Probabilistic}} {{Term}} {{Rewriting}} {{Using}}
               {{Dependency}} {{Pairs}}},
  booktitle = {Proc.\ CADE~'23},
  series    = {LNCS 14132},
  pages     = {344--364},
  year      = 2023,
  doi       = {10.1007/978-3-031-38499-8_20}
}

@inproceedings{ahrens2025WeightedRewritingSemiring,
  author    = {Ahrens, Emma and Kassing, Jan-Christoph and Giesl, J{\"u}rgen
               and Katoen, Joost-Pieter},
  title     = {{{Weighted}} {{Rewriting}}: {{Semiring}} {{Semantics}} for
               {{Abstract}} {{Reduction}} {{Systems}}},
  booktitle = {Proc.\ FSCD~'25},
  series    = {LIPIcs 337},
  pages     = {6:1--6:21},
  year      = 2025,
  doi       = {10.4230/LIPIcs.FSCD.2025.6}
}
\end{lstlisting}

Further venue examples: \lstinline|Proc.\ IJCAR~'24| (LNCS 14739),
\lstinline|Proc.\ TACAS~'26| (LNCS 16506), \lstinline|Proc.\ CAV~'26|
(LNCS 16684), \lstinline|Proc.\ FroCoS~'05| (LNCS 3717),
\lstinline|Proc.\ RTA~'10| (LIPIcs 6).
LNAI volumes are written as LNCS.
Joint events keep the joint acronym (\lstinline|Proc.\ RTA-TLCA~'14|).
Workshops without a series get only pages and year (\lstinline|Proc.\ WST~'23|,
\lstinline|Proc.\ IWC~'12|).

\paragraph{Journals.}
Use the ISO abbreviation, volume, pages, and year.
Add the issue number only when the pages do not identify the paper (LMCS,
article-numbered journals).
The following entries are rendered as \cite{fuhs2011provingterminationdependency},
\cite{arts2000termination}, and \cite{kassing2025LMCS}.

\begin{lstlisting}
@article{fuhs2011provingterminationdependency,
  author  = {Fuhs, Carsten and Giesl, J{\"u}rgen and Parting, Michael and
             {Schneider-Kamp}, Peter and Swiderski, Stephan},
  title   = {{{Proving}} {{Termination}} by {{Dependency}} {{Pairs}} and
             {{Inductive}} {{Theorem}} {{Proving}}},
  journal = {J.\ Autom.\ Reason.},
  volume  = 47,
  pages   = {133--160},
  year    = 2011,
  doi     = {10.1007/s10817-010-9215-9}
}

@article{arts2000termination,
  author  = {Arts, Thomas and Giesl, J{\"u}rgen},
  title   = {{{Termination}} of {{Term}} {{Rewriting}} {{Using}}
             {{Dependency}} {{Pairs}}},
  journal = {Theor.\ Comput.\ Sci.},
  volume  = 236,
  pages   = {133--178},
  year    = 2000,
  doi     = {10.1016/S0304-3975(99)00207-8}
}

@article{kassing2025LMCS,
  author  = {Kassing, Jan-Christoph and Giesl, J{\"u}rgen},
  title   = {{{From}} {{Innermost}} to {{Full}} {{Probabilistic}} {{Term}}
             {{Rewriting}}: {{Almost-Sure}} {{Termination}}, {{Complexity}},
             and {{Modularity}}},
  journal = {Log.\ Methods Comput.\ Sci.},
  volume  = 21,
  number  = 4,
  year    = 2025,
  doi     = {10.46298/lmcs-21(4:28)2025}
}
\end{lstlisting}

Other journal abbreviations:
\begin{lstlisting}
J.\ ACM                          Inf.\ Comput.
J.\ Symb.\ Comput.               Inf.\ Process.\ Lett.
ACM Trans.\ Comput.\ Log.        J.\ Log.\ Algebr.\ Methods Program.
\end{lstlisting}

\paragraph{Book chapters.}
Use \lstinline|@incollection| with \lstinline|booktitle|, \lstinline|editor|,
\lstinline|publisher|, series, and pages.
The following entry is rendered as \cite{brockschmidt2010TerminationGraphsJava}.

\begin{lstlisting}
@incollection{brockschmidt2010TerminationGraphsJava,
  author    = {Brockschmidt, Marc and Otto, Carsten and {von Essen}, Christian
               and Giesl, J{\"u}rgen},
  title     = {{{Termination}} {{Graphs}} for {{\tool{Java}}} {{Bytecode}}},
  booktitle = {Verification, {{Induction}}, {{Termination Analysis}}},
  editor    = {Siegler, Simon and Wasser, Nathan},
  publisher = {Springer},
  series    = {LNCS 6463},
  pages     = {17--37},
  year      = 2010,
  doi       = {10.1007/978-3-642-17172-7_2}
}
\end{lstlisting}

\paragraph{Theses.}
Use \lstinline|@phdthesis| or \lstinline|@mastersthesis| with the university
in \lstinline|school|.
The following entries are rendered as \cite{thiemann_diss:2007} and
\cite{kassing:thesis:2022}.

\begin{lstlisting}
@phdthesis{thiemann_diss:2007,
  author = {Thiemann, Ren{\'e}},
  title  = {{{The}} {{DP}} {{Framework}} for {{Proving}} {{Termination}} of
            {{Term}} {{Rewriting}}},
  school = {RWTH Aachen University},
  year   = 2007
}

@mastersthesis{kassing:thesis:2022,
  author = {Kassing, Jan-Christoph},
  title  = {{{Using}} {{Dependency}} {{Pairs}} for {{Proving}} {{Almost-Sure}}
            {{Termination}} of {{Probabilistic}} {{Term}} {{Rewriting}}},
  school = {RWTH Aachen University},
  year   = 2022
}
\end{lstlisting}

A bachelor's thesis is entered as \lstinline|@mastersthesis| with
\lstinline|type = {Bachelor's thesis}|.

\paragraph{Technical reports.}
Use \lstinline|@techreport| with \lstinline|number = {Technical Report ...}|
and \lstinline|institution = {...}|.

\paragraph{Online material.}
Use \lstinline|@misc| with the URL and access date in \lstinline|note|.
The following entry is rendered as \cite{aprove-website}.

\begin{lstlisting}
@misc{aprove-website,
  key  = {{\relax APrW}},
  note = {Website of \tool{AProVE}:
          \url{https://aprove.informatik.rwth-aachen.de/}.
          Accessed: 2026-07-19}
}
\end{lstlisting}

%-----------------------------------------------------------------------
\section{DOI}\label{sec:bib-doi}
%-----------------------------------------------------------------------

Every entry that has a DOI gets \lstinline|doi = {10.|\ldots\lstinline|}|.
The style renders it as a link.
Workshop papers without a DOI get the proceedings URL in \lstinline|note|.
Never put a URL into the \lstinline|doi| field.

%-----------------------------------------------------------------------
\section{Non-English Sources}\label{sec:bib-non-english}
%-----------------------------------------------------------------------

Add \lstinline|note = {In Japanese}| (or the respective language) for papers
in non-English journals.

%-----------------------------------------------------------------------
\section{Results}\label{sec:bib-results}
%-----------------------------------------------------------------------

With these fields the style produces, uniformly:

\begin{quote}
    Jan-Christoph Kassing and Jürgen Giesl. Proving Almost-Sure Innermost
    Termination of Probabilistic Term Rewriting Using Dependency Pairs. In
    \emph{Proc.\ CADE~'23}, LNCS 14132, pages 344--364, 2023.
    doi:10.1007/978-3-031-38499-8\_20.
\end{quote}

\begin{quote}
    Carsten Fuhs, Jürgen Giesl, Michael Parting, Peter Schneider-Kamp, and
    Stephan Swiderski. Proving Termination by Dependency Pairs and Inductive
    Theorem Proving. \emph{J.\ Autom.\ Reason.}, 47:133--160, 2011.
    doi:10.1007/s10817-010-9215-9.
\end{quote}

\begin{quote}
    Nils Lommen, Éléanore Meyer, and Jürgen Giesl. Control-Flow Refinement for
    Complexity Analysis of Probabilistic Programs in \tool{KoAT} (Short Paper).
    In \emph{Proc.\ IJCAR~'24}, LNCS 14739, pages 233--243, 2024.
    doi:10.1007/978-3-031-63498-7\_14.
\end{quote}

These are the entries \cite{kassinggiesl2023iast},
\cite{fuhs2011provingterminationdependency}, and \cite{lommen2024KoATShortPaper}
in the bibliography of this thesis.

%-----------------------------------------------------------------------
\section{Ordering}\label{sec:bib-ordering}
%-----------------------------------------------------------------------

The bibliography is sorted alphabetically by author, not by order of citation:
first by the surname of the first author, then by the following authors, then
by year, then by title.
Entries without authors (websites, tools) sort under their \lstinline|key|
field, which also determines their label (see the \lstinline|@misc| example
in \cref{sec:bib-venues}).
Particles are part of the surname for sorting, so
\lstinline|{von Essen}, Christian| is filed under V (between Thiemann and
Wasser), not under E.
For the examples above this gives the order
\begin{quote}
    Arts, Giesl (2000) -- Brockschmidt, Otto, von Essen, Giesl (2010) --
    Fuhs, Giesl, Parting, Schneider-Kamp, Swiderski (2011) --
    Kassing, Giesl (2023) -- Kassing, Giesl (2025) --
    Lommen, Meyer, Giesl (2024) -- Thiemann (2007).
\end{quote}

Keep the \lstinline|.bib| file itself in the same order (e.g., with
\lstinline|bibtool -s|, sort key first author and year), so that
inconsistencies between entries of the same author are easy to spot.
